\documentclass{ximera}

\input{../../preamble.tex}

\title{Elementary Library 1}

\begin{document}

\begin{abstract}
%Stuff can go here later if we want!
\end{abstract}
\maketitle


The \textbf{Core Functions} is a collection of basic building blocks for functions.

We build other functions by combining the Core Functions forming sums, difference, products, quotients, and compostions. 

The collection of all such combinations is called The \textbf{Elementary Functions}.

The \textbf{Elementary Functions} includes some very familiar functions.

They include 

\begin{itemize}
\item constant, linear, quadratic, cubic, and other polynomial functions
\item rational functions
\item root or radical functions
\item exponential and logarithmic functions
\item trigonometric functions
\item absolute value, greatest integer, other individual piecewise defined functions
\end{itemize}



This section and then next will just introduce the library of elementary functions with some basic defining features.  This section is the beginning of their analysis.












\subsection*{Learning Outcomes}

\begin{sectionOutcomes}
In this section, students will 

\begin{itemize}
\item be introduced to the library.
\item identify basic defining features and characteristics.
\end{itemize}
\end{sectionOutcomes}







\begin{center}
\textbf{\textcolor{green!50!black}{ooooo-=-=-=-ooOoo-=-=-=-ooooo}} \\

more examples can be found by following this link\\ \link[More Examples of Elementary Functions]{https://ximera.osu.edu/csccmathematics/precalculus1/precalculus1/elementaryLibrary1/examples/exampleList}

\end{center}




\end{document}

